% sec-01-abstract.tex - Abstract and Keywords (full)
\begin{abstract}
\noindent
A Phase 1 oncology trial advances only when a sequence of large regulatory and
clinical documents is authored, reviewed, and approved, while patients wait for that paperwork to be assembled. This paper illustrates a
repository based large language model (LLM), driven by a single master prompt that writes and executes a schedule of sub-prompts committed to GitHub in real time; which can hasten the entire
Phase 1 process by generating documents through an auditable
"mermaid" to "draft" to "full" to "final" pipeline. The method was conducted on a prior
pancreatic ductal adenocarcinoma program: an on-premises LLM directed
robotic Whipple with perioperative daraxonrasib (RMC-6236), whose Phase 1 and
Phase 2 protocols were themselves single-prompt multi-file builds. The following trial aspects were addressed: the initial IND and IRB package; protocol amendments with synchronized
consent; cohort-review packages; the complete clinical-hold response; the
Phase 2-to-3 briefing package; and the CSR and NDA/BLA modules. Only the administrative and preparation time bucket compresses, while
clinical follow-up and fixed regulatory review clocks do not. Papers finish in 1-4 days, based on prompt and input status. Twenty-four colored Mermaid converted figures ground the prose in
LaTeX. Paper repository, directory, and file name samples were verified manually by the author. For terminal patients, probable benefit exceeds probable risk.  
\end{abstract}

\keywords{Pancreatic ductal adenocarcinoma; Phase 1 clinical trial; document generation; repository based LLM; sub-prompt schedule; clinical trial acceleration; daraxonrasib.}
